% Problemas del capítulo 8. Generado por tools/problemas/gen.py; no editar a mano.
\section{Problemas}\label{sec:problemas-08}

Graduados: \({\star}\) directo, \({\star}{\star}\) para pensar, \({\star}{\star}{\star}\) desafío. Las respuestas están en el apéndice~\ref{ap:respuestas}; conviene intentarlos antes de mirarlas.

\begin{problema}[\({\star}\)]\label{pb:08-1}
¿En qué elipse convierte \(A=\begin{pmatrix}3&0\\0&1\end{pmatrix}\) a la circunferencia unidad? ¿Cuáles son sus valores singulares?
\end{problema}

\begin{problema}[\({\star}{\star}\)]\label{pb:08-2}
Hacé lo mismo con \(A=\begin{pmatrix}2&1\\1&2\end{pmatrix}\). ¿En qué direcciones quedan los ejes de la elipse?
\end{problema}

\begin{problema}[\({\star}{\star}\)]\label{pb:08-3}
Calculá los valores singulares de la cizalla \(A=\begin{pmatrix}1&1\\0&1\end{pmatrix}\) y comprobá que su producto es \(|\det A|\).
\end{problema}

\begin{problema}[\({\star}{\star}\)]\label{pb:08-4}
La incertidumbre de un punto medido tiene matriz de covarianza \(\Sigma_{EN}=\begin{pmatrix}4&1{,}5\\1{,}5&1\end{pmatrix}\)~cm\(^2\), en \((E,N)\). Calculá los semiejes de la elipse de error (las raíces de los autovalores) y el acimut del eje mayor.
\end{problema}

\begin{problema}[\({\star}{\star}{\star}\)]\label{pb:08-5}
Para \(A=\begin{pmatrix}1&1\\1&1{,}01\end{pmatrix}\), calculá los valores singulares y su cociente. ¿Qué forma tiene la elipse? ¿Qué pasa si resolvés \(Ax=b\) con datos que tienen un poco de error?
\end{problema}
